% ============================================================================
% Chapter 3 — System Overview
% ============================================================================
\chapter{System Overview}
\label{ch:overview}

This chapter provides a high-level walkthrough of the \versimux{} system before the detailed architectural exposition in \cref{ch:architecture}. It introduces the design goals, presents the end-to-end pipeline, describes the technology stack, and discusses key architectural trade-offs.


% ---- 3.1 Design Goals and Principles ----------------------------------------
\section{Design Goals and Principles}
\label{sec:design-goals}

The design of \versimux{} is guided by five principles:

\begin{description}
  \item[Grounding.] Every agent action must be validated against the live \acrfull{DOM}{Document Object Model} before execution. Selectors that do not exist in the current page state are rejected, preventing hallucinated interactions.

  \item[Closed-Loop Remediation.] The system does not terminate at diagnosis. It synthesises concrete, typed code patches (HTML, CSS, JavaScript), applies them, and verifies their effectiveness---re-entering the evaluation loop if the resolution threshold is not met.

  \item[Actionability.] Output artifacts must be directly usable by developers: code patches that can be applied to source files, not textual suggestions that require manual interpretation.

  \item[Observability.] Every pipeline stage emits structured events that are persisted, streamed to the frontend in real time, and available for post-hoc analysis.

  \item[Provider Agnosticism.] The \acr{LLM} integration layer abstracts over multiple providers (Groq, OpenAI, DeepSeek, Qwen, Moonshot), enabling model substitution without architectural changes.
\end{description}


% ---- 3.2 End-to-End Pipeline Walkthrough ------------------------------------
\section{End-to-End Pipeline Walkthrough}
\label{sec:pipeline-walkthrough}

\begin{figure}[htbp]
  \centering
  \framebox[\textwidth]{\rule{0pt}{7cm}\parbox{0.9\textwidth}{\centering\large\sffamily [Figure 3.1: Full System Architecture Diagram]\\ \vspace{1em} \normalsize\normalfont Displays the three-tier design of \versimux{}: the Next.js Nexus client, the FastAPI backend API server with SQLite and the event bus, and the core LangGraph multi-agent pipeline.}}
  \caption{Overall system architecture showing the three-tier design: frontend client, backend API server, and the LangGraph agent pipeline.}
  \label{fig:system-architecture}
\end{figure}

The \versimux{} pipeline processes a single HTML document through twelve sequential stages:

\begin{enumerate}
  \item \textbf{HTML Upload.} The user submits an HTML file and an optional UI context description via the REST \acrfull{API}{Application Programming Interface}.

  \item \textbf{Supervisor Analysis.} The supervisor agent analyses the HTML structure, identifies interactive elements, critical user paths, and accessibility risk indicators.

  \item \textbf{Persona Generation.} Based on the analysis, the supervisor generates a diverse set of persona profiles with distinct accessibility constraints, cognitive characteristics, and task goals.

  \item \textbf{Persona Simulation.} Each persona agent executes its task goal inside an isolated Playwright browser context, producing an interaction trace with screenshots and inline issue annotations.

  \item \textbf{Trace Verification.} Rule-based heuristics validate each trace for structural integrity: loop detection, selector existence, navigate interception, and action validity.

  \item \textbf{Issue Extraction.} Verified issues from all personas are aggregated into a unified issue set.

  \item \textbf{Issue Clustering.} Sentence-transformer embeddings and HDBSCAN clustering group semantically similar issues, eliminating cross-persona duplicates.

  \item \textbf{Recommender Profile Generation.} The supervisor creates a \code{RecommenderProfile} for each cluster, specifying the fix focus, strategy hint, and priority.

  \item \textbf{Patch Synthesis.} Parallel recommender agents propose typed patches for their assigned clusters.

  \item \textbf{Conflict Resolution.} A conflict resolver detects overlapping patches and resolves them through LLM-mediated negotiation or heuristic tiebreaking.

  \item \textbf{Patch Application \& Verification.} Patches are applied to the HTML in dependency order. A verification agent assesses whether targeted issues have been resolved. If the resolution threshold is not met, the pipeline enters a correction loop.

  \item \textbf{Report Generation.} A comprehensive diagnostic report is generated with scores, issue details, applied patches, and verification results.
\end{enumerate}

\begin{figure}[htbp]
  \centering
  \framebox[\textwidth]{\rule{0pt}{7cm}\parbox{0.9\textwidth}{\centering\large\sffamily [Figure 3.2: Pipeline Stage Sequence Diagram]\\ \vspace{1em} \normalsize\normalfont Illustrates the temporal progression and state changes across the twelve pipeline stages, highlighting the iterative correction loop between patch verification and persona re-simulation.}}
  \caption{The twelve-stage pipeline sequence from HTML upload to report generation.}
  \label{fig:pipeline-sequence}
\end{figure}


% ---- 3.3 Technology Stack ----------------------------------------------------
\section{Technology Stack}
\label{sec:tech-stack}

\Cref{tab:tech-stack} summarises the key technologies employed at each architectural layer.

\begin{table}[htbp]
  \centering
  \caption{Technology stack by architectural layer.}
  \label{tab:tech-stack}
  \begin{tabularx}{\textwidth}{@{}l L{3cm} X@{}}
    \toprule
    \textbf{Layer} & \textbf{Technology} & \textbf{Purpose} \\
    \midrule
    Orchestration      & LangGraph            & State-graph pipeline with typed state, conditional edges, and parallel fan-out via \code{Send()} \\
    \acr{LLM} Backend  & Groq, OpenAI, etc.   & Provider-agnostic \acr{LLM} calls with per-role model routing and rate limiting \\
    Browser Automation & Playwright (Chromium) & Headless browser control with isolated contexts per persona \\
    Embeddings         & sentence-transformers & Issue text embedding for semantic clustering \\
    Clustering         & HDBSCAN              & Density-based clustering with automatic cluster count \\
    Backend \acr{API}  & FastAPI              & REST endpoints and \acr{SSE} streaming \\
    Persistence        & SQLite               & Session and event storage \\
    Frontend           & Next.js (React)      & Dashboard, upload interface, and report viewer \\
    \bottomrule
  \end{tabularx}
\end{table}


% ---- 3.4 Architectural Constraints and Trade-offs ----------------------------
\section{Architectural Constraints and Trade-offs}
\label{sec:trade-offs}

Several architectural decisions reflect deliberate trade-offs:

\paragraph{LangGraph over linear orchestration.}
Linear orchestration scripts are inherently brittle when applied to complex agentic workflows that require branching, error recovery, or feedback loops. A simple sequential execution path cannot easily accommodate the loopback flows required for closed-loop remediation, where failing verification results trigger a return to persona simulation. By using LangGraph, \versimux{} models the evaluation and remediation pipeline as a directed state graph. This formal structure supports conditional edges to route execution back to earlier phases, parallel fan-out via the \code{Send()} primitive to evaluate multiple HTML documents and personas concurrently, and typed state accumulation. Crucially, LangGraph's state design allows parallel branches to merge their results into a unified, type-safe schema using custom reducer functions (e.g., \code{operator.add}), eliminating the need to write custom concurrency locks or synchronization primitives.

\paragraph{ThreadPoolExecutor over asyncio for persona simulation.}
While asynchronous I/O via Python's \code{asyncio} is highly efficient for network-bound operations, integrating it with browser automation libraries like Playwright introduces significant event-loop complexity when scaled across multiple concurrent instances. Parallel browser automation runs frequently block the event loop, causing timeouts and race conditions in LLM prompt generation and DOM interaction. Furthermore, if a single page simulation crashes or hangs, isolating the event loop failure can be challenging. \versimux{} resolves these reliability concerns by using sync-mode Playwright wrapped in OS-level threads managed by a \code{ThreadPoolExecutor}. This model provides robust thread-local isolation: each persona agent maintains its own thread, its own Chromium browser context, and independent error boundaries. The slightly higher memory overhead of system threads is a worthwhile trade-off for the increased stability and cleaner exception handling during concurrent simulations.

\paragraph{SQLite over Redis for event persistence.}
To support real-time user-interface telemetry and SSE streaming, the backend requires a persistent database to record job events, telemetry logs, and diagnostic states. While dedicated in-memory stores like Redis are standard for high-throughput streaming architectures, they introduce external system daemon dependencies, increasing the complexity of setup and local deployment. Because \versimux{} is designed as a developer-facing tool running locally or in single-host container environments, high-throughput clustering is unnecessary. A local SQLite database configured with Write-Ahead Logging (WAL) and synchronous mode set to normal offers sufficient write performance (exceeding hundreds of transactions per second) and concurrent read safety, while maintaining a single-file, zero-dependency deployment model.

\paragraph{Python-maintained working memory over LLM self-tracking.}
In multi-turn web navigation simulations, the persona agent must track its operational state (e.g., active page context, target form fields filled, remaining interaction steps, and past action history). Allowing the underlying \acr{LLM} to self-track this state through prompt instructions or conversational memory inevitably leads to cognitive drift, where the model forgets its remaining step budget, hallucinates having filled non-existent inputs, or enters repetitive interaction loops. \versimux{} avoids this failure mode by keeping the agent's core working memory strictly in Python data structures (specifically, the \code{WorkingMemory} dataclass). After each browser interaction, the Python runtime updates the state variables, counts down the step budget, and formats this structured state to be injected into the next LLM prompt. This hybrid approach ensures that the model remains grounded in physical reality, preventing context hallucination while leveraging the LLM's reasoning power for action decision-making.
